\section{Experimental Evaluation}
\label{sec:experiments}

Our experiments ask four questions. Does the search return plans with the
minimum number of grasps for new goals, start configurations and obstacle
layouts, from a single offline stage? How does it compare with a mixed-integer
formulation and with a sampling-based manipulation planner as the number of
required grasps grows? Why must the transition sets be validated? And does validation
change the plans?

\subsection{Tasks and implementation}
Both tasks turn a cap about a fixed vertical axis with a Franka
Panda~\cite{haddadin2022franka}, whose wrist is restricted to $[-1,1]$\,rad,
so any turn larger than 2\,rad needs several grasps.
\begin{itemize}
  \item \textbf{T1 (3-DOF).} All arm joints but the wrist are locked:
    $q=(w,x_{\mathrm{finger}},\alpha)$, with
    $x_{\mathrm{finger}}\in[-0.045,0]$\,m and $\alpha\in[-\pi,\pi]$. The
    decomposition has 15 regions, and the task turns the cap by $2\pi$.
    \tocheck{T1 from the MInf report (WSG gripper) or the locked Franka?}
  \item \textbf{T2 (10-DOF).} Seven arm joints, two fingers and the cap. The
    cap is 0.5\,m in front of the robot base, with two cylindrical obstacles
    (radius 5\,mm) next to it. The free sets come from an IRIS-NP2 clique
    cover with 42 regions \tocheck{coverage target 0.98; $(\varepsilon,\delta)$;
    generation time}. The reference query Q0 turns the cap from $-1$ to 3\,rad
    and returns the arm to its start configuration.
\end{itemize}

The grasp constraints~\eqref{eq:grasp_hg} ask the hand to be centred over the
cap, within a height band, with its approach axis $z_E={}^{C}R^{E}(x)e_z$
pointing down:
\begin{align}
  h(x) &= [\,p_x,\; p_y,\; z_{E,x},\; z_{E,y}\,]^\top, \label{eq:grasp_h}\\
  g(x) &= [\,p_z - z_{\max},\; z_{\min} - p_z,\; z_{E,z} + c\,]^\top, \label{eq:grasp_g}
\end{align}
with $p={}^{C}p^{E}(x)$, $[z_{\min},z_{\max}]=[0.105,0.110]$\,m and $c=0.9$,
and the fingers held within a 1\,mm band around the width of the cap. Here
$J_h$ has rank 4, leaving a six-dimensional null space. The grasp set
of~\eqref{eq:graspPoly} uses $\epsilon=10^{-3}$, $\rho=0.05$\,rad and
$\Lmin=0.05$\,rad. The wrist of the Panda satisfies the invariance of
Section~\ref{sec:method:grasp} exactly: sweeping it over its whole range at a
grasp leaves $h$ and $g$ unchanged to machine precision.
The planner is implemented in Python with Drake~\cite{tedrake2019drake}
(v1.48). The LPs use Drake's \texttt{MathematicalProgram}, and the MIQP uses
MOSEK~\cite{mosek}. \todo{hardware; code release.}

\subsection{Baselines and protocol}
\label{sec:experiments:protocol}
We compare against three baselines:
\begin{itemize}
  \item GCS~\cite{marcucci2023motion} on the same sets and the same graph as
    the search (below);
  \item the mixed-integer formulation of Section~\ref{sec:method:miqp}, with
    GCS connecting its keyframes;
  \item a sampling-based planner on the manipulation
    graph~\cite{simeon2004manipulation}, with probabilistic roadmaps for transit
    and transfer. \todo{implement the sampling-based baseline (own
    implementation or HPP~\cite{mirabel2016hpp}); time limit.}
\end{itemize}
For all methods, every returned motion is checked densely with Drake's
collision checker.

\paragraph{GCS on the same sets}
The free sets $F_j$ and the transition sets $T_k$, with $\alpha$ free within its
limits, are the vertices, and the transit and switch edges of
Section~\ref{sec:method:graph} are the edges. Each visited vertex holds one
straight segment $r_0\to r_1$, with both ends in its set. The coupling
becomes linear constraints on these two points: $r_1^{\alpha}=r_0^{\alpha}$ in
every $F_j$, and in every $T_k$
\[
r_1^{\alpha}-r_1^{w}=r_0^{\alpha}-r_0^{w}, \qquad r_1^{w}\ge r_0^{w},
\]
which is a forward stroke. A GCS path visits each vertex at most once, but a
plan can need several strokes in the same $T_k$, with returns in between. So
the graph is copied once per grasp:
\[
\text{approach}\to\text{grasp}_1\to\text{between}_1\to\cdots\to\text{grasp}_N\to\text{depart},
\]
with $N=\lceil D/\Lmax\rceil+1$ for a turn $D$. Every $\text{grasp}_n$ also
connects to the depart layer, and the approach layer also connects directly to
the goal.

The cost is the path length, plus 10 per grasp (each free-to-transition edge),
plus $\mu$ per other change of set. Ten exceeds any path length on T2, so, as in
the search, the fewest grasps come first. The between layers hold only the free
sets that meet a transition set. The approach and depart layers hold those,
the free sets containing $\xinit$ or $\xgoal$, and the free sets used by the A*
plan. The reason for this pruning is memory (Section~\ref{sec:results:gcs}).

We solve this graph with Drake's \texttt{SolvePath}: a convex relaxation,
rounded to at most 10 candidate paths, each finished by a convex restriction,
with MOSEK. We also solve only the convex restriction on the sequence of sets
that the A* returns (\emph{GCS on the A* sequence}). It keeps the A*'s discrete
plan and optimises its configurations, so it isolates the second stage
(Section~\ref{sec:method:recovery}).

Five experiments follow. \textbf{E1} sweeps $\agoal$ over the cap range at a
fixed arm configuration, and reports the number of grasps against the lower
bound $\lceil D/\Lmax\rceil$ for a turn $D=\agoal-\ainit$. \textbf{E2} draws random start and
goal configurations from the free sets. \textbf{E3} varies the obstacle
layout, each layout with its own offline stage, and separates offline from
online cost. \textbf{E4} compares the reachability A* with the baselines as a function of
the number of grasps a query requires, reporting success within the time limit,
planning time and grasps. \textbf{E5} runs Q0 with and without the validation
of Section~\ref{sec:method:graph}.
\todo{run E1--E5; this is the evidence for the main claim.}

\section{Results}
\label{sec:results}

\subsection{Reachability A* on the 10-DOF task}
Table~\ref{tab:franka} and Fig.~\ref{fig:reach} report Q0, before validation.
Of the 42 free sets, 8 intersect the grasp set. Four of these transition sets
span the whole wrist range ($L_k=2$\,rad), and the others span between 0.06 and
0.36\,rad. The search expands 63 nodes in 0.004\,s and returns two strokes, one
return and one release: it enters $T_6$ from $F_2$, performs stroke, return
and stroke, and releases into $F_0$. Its cost, $4+\mu$, equals the initial
heuristic plus one set change, so the heuristic is tight on this query. With
$D=4$\,rad and $\Lmax=2$\,rad, two grasps is the minimum that the sets allow.
The second stage turns 2\,rad per stroke and produces 7 waypoints, and all 6
straight-line motions lie in their recorded set and satisfy the coupling. The online stage performs no
geometric computation of its own: every intersection and extreme point it
relies on was computed offline, by linear programs.

\begin{table}[t]
\centering
\caption{Reachability A* on T2 (10-DOF), query Q0, before validation}
\label{tab:franka}
\setlength{\tabcolsep}{4pt}
\begin{tabular}{@{}l l r@{}}
\toprule
Stage & Quantity & Value\\
\midrule
\multirow{5}{*}{Offline} & free sets $J$ (non-empty after slicing) & 42\\
 & transition sets $K$ (with $L_k=2$\,rad) & 8 (4)\\
 & transit / switch edges & 850 / 58\\
 & graph construction time & \tbd\\
 & decomposition time & \tbd\\
\midrule
\multirow{5}{*}{Online} & nodes expanded & 63\\
 & search time & 0.004\,s\\
 & strokes / returns / releases & 2 / 1 / 1\\
 & set changes & 1\\
 & motions contained in their set & 6 / 6\\
\bottomrule
\end{tabular}
\end{table}

\subsection{The trust region on the grasp set}
Without the trust region $\rho$, the extreme points that the graph
construction returns lie far outside the region where the linearisation holds:
on T2 they lay 2.27\,rad from the seed in one arm joint, with
$\|h\|_\infty=0.95$ and the hand 0.80\,m away from the cap. With
$\rho=0.05$\,rad, $\|h\|_\infty$ stays near $2\times10^{-3}$ along the plans we
report \tocheck{numbers from the q0 run}.

\subsection{Comparison with the MIQP}
Table~\ref{tab:miqp} compares the two ways of choosing grasps. On T1, the MIQP
returns 11 grasps after proving 10 infeasible. Each stroke turns the cap by
0.604\,rad, the wrist range of the grasp set in that decomposition, and
$\lceil 2\pi/0.604\rceil=11$. GCS then connects all 23 transit and transfer
segments. On T2 it does not return a reliable solution. The
numbers make the argument of Section~\ref{sec:method:miqp} concrete: the grasp
set is $2\epsilon=2\times10^{-3}$ thick, while $M=10^{6}$ and MOSEK's default
integrality tolerance $\tau=10^{-5}$~\cite{mosek} give $M\tau=10$\,rad of
slack on a set reported as occupied, more than any joint range. The
reachability A* returns the plan above in 0.004\,s. \tocheck{MIQP outcome and solve time on T2; run the
reachability A* on T1 and check that it also returns 11.}

\begin{table}[t]
\centering
\caption{Choosing grasps with the MIQP and with the reachability A*}
\label{tab:miqp}
\setlength{\tabcolsep}{4pt}
\begin{tabular}{@{}l l c c l@{}}
\toprule
Task & Method & Grasps & Time & Outcome\\
\midrule
\multirow{2}{*}{T1 (3-DOF)} & MIQP & 11 & \tbd & solved; GCS connected\\
 & Reach.\ A* & \tbd & \tbd & \tbd\\
\midrule
\multirow{2}{*}{T2 (10-DOF)} & MIQP & -- & \tbd & numerical failure \tocheck{}\\
 & Reach.\ A* & 2 & 0.004\,s & solved (rel.\ sets)\\
\bottomrule
\end{tabular}
\end{table}

\subsection{Comparison with GCS on the same sets}
\label{sec:results:gcs}

\paragraph{Setup}
\begin{itemize}
  \item All runs use T2 with the 42 free sets. The transition sets are
    restricted to the wrist intervals that the collision checker accepts, so
    $\Lmax=0.81$\,rad. This was implemented by cutting the free sets with planes
    normal to the wrist where the wrist sweep of
    Section~\ref{sec:results:fidelity} finds collisions.
    \tocheck{same as the validation of Section~\ref{sec:method:graph}?}
  \item Every row of Table~\ref{tab:gcs} shares one start and one goal, saved
    with the sets so that the run reproduces exactly. The cap starts at
    $\ainit=-3$\,rad and turns by $D=0.5$, 1.5 or 6\,rad, which needs 1, 2 or
    8 grasps. The code uses the opposite sign.
  \item Times count each method's online work:
    \begin{itemize}
      \item reachability A*: the search and the placement of the waypoints;
      \item GCS on the A* sequence: the A* search, building the layered GCS
        graph, and one convex program;
      \item full GCS: building the layered graph and \texttt{SolvePath}
        (relaxation, then rounding to at most 10 candidate paths).
    \end{itemize}
  \item Not counted for any method: the 1.1\,s that builds the A*'s graph,
    which GCS reuses instead of repeating its intersection checks; and the
    decomposition.
  \item Memory is the peak resident memory of the computation. Each full GCS
    solve runs in its own process. MOSEK solves all programs.
  \item Appendix~\ref{app:gcs} lists every run, including those that failed.
\end{itemize}
Before these runs we checked the encoding on a two-variable toy, not a scene:
$(w,\alpha)$ in a box, a 4\,rad turn, strokes of at most 2\,rad. GCS returned
stroke, return and stroke, as the search does. With a single copy of the grasp
vertex it is infeasible, as the next paragraph explains.

\begin{table}[t]
\centering
\caption{Plans on T2 from one start and goal. Length: Euclidean, in
$(x,\alpha)$. Every plan is certified: each motion lies in its convex set and
respects the coupling.}
\label{tab:gcs}
\setlength{\tabcolsep}{3.5pt}
\begin{tabular}{@{}r l r r r r@{}}
\toprule
$D$ [rad] & Method & Grasps & Length & Time [s] & Mem. [GB]\\
\midrule
\multirow{3}{*}{0.5} & Reach.\ A* & 1 & 18.06 & 0.0037 & $<$0.001\\
 & GCS on A* seq. & 1 & 11.49 & 0.44 & $<$0.001\\
 & Full GCS & 1 & 11.49 & 191 & 3.0\\
\midrule
\multirow{3}{*}{1.5} & Reach.\ A* & 2 & 20.30 & 0.0053 & $<$0.001\\
 & GCS on A* seq. & 2 & 13.62 & 0.54 & 0.007\\
 & Full GCS & 2 & 13.62 & 403 & 4.3\\
\midrule
\multirow{3}{*}{6} & Reach.\ A* & 8 & 31.16 & 0.014 & $<$0.001\\
 & GCS on A* seq. & 8 & 24.48 & 0.54 & 0.025\\
 & Full GCS & \multicolumn{4}{l}{did not fit in 60\,GB (Table~\ref{tab:gcs:scaling})}\\
\bottomrule
\end{tabular}
\end{table}

\begin{table}[t]
\centering
\caption{Cost of each method against the number of grasps. The A*'s graph is
the same for every query: 50 vertices (42 free and 8 transition sets) and 900
edges. The GCS graph is copied once per grasp. Memory: GB.}
\label{tab:gcs:scaling}
\setlength{\tabcolsep}{3pt}
\begin{tabular}{@{}r r r r r r@{}}
\toprule
 & \multicolumn{2}{c}{Reach.\ A*} & \multicolumn{3}{c}{Full GCS}\\
\cmidrule(lr){2-3}\cmidrule(l){4-6}
Grasps & Nodes & Time [s] & $|V|$\,/\,$|E|$ & Time [s] & Mem.\\
\midrule
1 & 83 & 0.0037 & 48 / 543 & 191 & 3.0\\
2 & 122 & 0.0053 & 64 / 749 & 403 & 4.3\\
3$^{a}$ & -- & 0.006 & 95 / 1461 & 680 & $\le$18.2\\
4$^{a}$ & -- & 0.005 & 111 / 1667 & 200 & $\le$18.2\\
8 & 369 & 0.014 & 179 / 2280$^{b}$ & -- & $>$15$^{b}$\\
\bottomrule
\end{tabular}

\vspace{2pt}
\footnotesize
\begin{minipage}{\columnwidth}
$^{a}$~From another run (Run~D in Appendix~\ref{app:gcs}), with its own
start and goal; memory there is an upper bound. $^{b}$~From another run (Run~E), with its own start and goal. The
relaxation solved in 140\,s at 13.7\,GB; the rounding was stopped at our
15\,GB limit. Without
pruning, the 8-grasp graph has about 20k edges and ran out of the machine's
60\,GB.
\end{minipage}
\end{table}

\paragraph{Results}
\begin{itemize}
  \item All methods return the same number of grasps, $\lceil D/\Lmax\rceil$.
  \item GCS on the A* sequence keeps the A*'s sets and grasps, and shortens its
    path by 21--36\%. Almost all of its 0.44--0.54\,s is spent building the
    layered graph; the convex program itself takes 0.011--0.028\,s.
  \item The A*'s paths are longer because its second stage places waypoints
    without optimising length (Section~\ref{sec:method:recovery}). That convex
    program is a natural replacement for the second stage.
  \item The full GCS returns the same path as GCS on the A* sequence, to four
    digits. Searching over all sets picked the sets the A* had already picked,
    at 191--403\,s and 3--4\,GB, against 0.004--0.005\,s for the search.
  \item Its cost grows with the number of grasps
    (Table~\ref{tab:gcs:scaling}), and at 8 grasps it did not fit on the
    machine. The A*'s cost barely changes: 83 to 369 nodes, 0.004 to 0.014\,s.
\end{itemize}

\paragraph{Why GCS needs the graph unrolled}
\begin{itemize}
  \item \emph{The rule.} In GCS a path is a sequence of distinct vertices, and
    each vertex holds a single point~\cite{marcucci2024shortest}. A regrasp
    plan, however, comes back to the same transition set once per grasp, and
    to the free sets where it lets go and turns the wrist back.
  \item \emph{The workaround.} GCS can represent such a plan only after the
    graph is unrolled by hand. We used one copy of the transition sets and of
    the free sets used between strokes for each grasp, an approach and a
    depart layer, the coupling written as linear constraints on each copy,
    and the number of grasps bounded before solving. With a single copy, the
    two-grasp toy above is infeasible, both relaxed and as the exact
    mixed-integer program.
  \item \emph{The cost.} The convex relaxation keeps both endpoints of the
    segments for every edge, so its size grows with the number of grasps times
    the edges of a copy. On T2, 850 of the 861 pairs of free sets intersect, so
    a copy of all of them adds about 1700 edges, and we had to prune the
    copies to fit.
  \item \emph{The reachability A* needs none of this.} A node pairs a set with
    an interval of $\alpha$, so coming back to a set with a different $\alpha$
    is simply a new node. The number of grasps is an output of the search, not
    a bound given in advance, and the graph is built once for all queries.
\end{itemize}
So the comparison does not show that GCS cannot plan regrasps, but that it
needs problem-specific engineering to do so, and that its cost then grows
with the number of grasps that the A* handles natively.
\todo{Full GCS with returns inside copies of the transition sets instead of
free sets (the A*'s model), which would roughly halve the graph; rerun the
8-grasp query with a higher memory limit.}

\subsection{Why transition sets must be validated}
\label{sec:results:fidelity}
A planner whose guarantees are relative to its sets inherits their errors. To measure them where they matter, we swept the wrist over $[-1,1]$
at the seed grasp, at 99 points, and compared membership in the free sets
with Drake's collision checker. The checker finds the wrist free only on
$[-1,-0.19]$ and $[0.70,0.92]$\,rad \tocheck{which bodies collide in between:
fingers vs.\ obstacles?}. All four decompositions we computed for T2
(IRIS-NP2 with 42 and 46 regions, and two IRIS-ZO covers with 64 regions
\tocheck{the two IRIS-ZO files may be the same cover}) contain every point of
the sweep, so all four label as collision-free the 46 configurations that the
checker rejects, and miss none of the 53 it accepts. The two-grasp plan of Fig.~\ref{fig:reach} is therefore valid
relative to the sets, but collides in reality. With the longest truly free
wrist interval of 0.81\,rad, the minimum is $\lceil4/0.81\rceil=5$ grasps: a
different discrete structure, not just a different geometry.

This is not a bug in the region generators: their guarantee bounds the
fraction of a region's volume in collision, and transition sets have almost
no volume (Section~\ref{sec:method:graph}). We observed the same sensitivity on low-dimensional versions of the task,
where a decomposition at 0.90 coverage contained a grasp and release sequence
that GCS could not connect. For regrasp planning, what matters is not coverage of free
space but the fidelity of the transition sets.
\todo{E5: Q0 with validated transition sets; expected 5 grasps with no
change to the search.}

\subsection{Generalisation across goals, starts and scenes}
\todo{Results of E1--E4 (Section~\ref{sec:experiments:protocol}).}

% Results figure, to add once E1/E4 are run: (left) grasps vs. turning
% angle D, a staircase tracking ceil(D/L); (right) planning time vs.
% required grasps, for the reachability A*, the MIQP and the
% sampling-based baseline. Uncomment and provide figures/generalisation.pdf.
% \begin{figure}[t]\centering
%   \includegraphics[width=\linewidth]{figures/generalisation.pdf}
%   \caption{...}\label{fig:generalisation}
% \end{figure}
